% ============ 第1章 §1.13 BMO空间(原书 pp.114-122) ============
\section{BMO 空间}\label{sec:1.13}
\index{BMO 函数}
\index{有界平均振动函数}
\index{John--Nirenberg 定理}


作为本章的结束,我们讨论 Hardy 空间 $H^{1}(\mathbb{R})$ 的对偶空间,同时也是与锐极大函数有关的函数空间----BMO 空间.

\begin{definition}\label{def:1.13.1}
若 $f \in L_{\mathrm{loc}}^{1}(\mathbb{R})$ 满足
\begin{equation*}
\| f \|_{\mathrm{BMO}} := \sup_{Q}\frac{1}{|Q|}\int_{Q}|f(x) - f_{Q}|\dd x < +\infty,
\end{equation*}
则称 $f$ 为 BMO 函数(有界平均振动函数).这里,上确界对一切区间 $Q$ 取,$f_{Q} = |Q|^{-1}\int_{Q}f(x)\dd x$.
\end{definition}

容易看出,$\| f \|_{\mathrm{BMO}} = \sup_{x \in \mathbb{R}}f^{\#}(x)$.

\begin{lemma}\label{lem:1.13.2}
对 $f \in L_{\mathrm{loc}}^{1}(\mathbb{R})$,
\begin{equation*}
\sup_{Q}\inf_{c \in \mathbb{C}}\frac{1}{|Q|}\int_{Q}|f(x) - c|\dd x \leq \sup_{Q}\frac{1}{|Q|}\int_{Q}|f(x) - f_{Q}|\dd x
\end{equation*}
\begin{equation*}
\leq 2\sup_{Q}\inf_{c \in \mathbb{C}}\frac{1}{|Q|}\int_{Q}|f(x) - c|\dd x.
\end{equation*}
\end{lemma}

\begin{proof}
对区间 $Q$ 与任意 $c \in \mathbb{C}$,
\begin{align*}
\dfrac{1}{|Q|}\displaystyle\int_{Q}|f(x) - f_{Q}|\dd x & \leq \dfrac{1}{|Q|}\int_{Q}|f(x) - c|\dd x + \dfrac{1}{|Q|}\int_{Q}|c - f_{Q}|\dd x \\
& = 2\dfrac{1}{|Q|}\displaystyle\int_{Q}|f(x) - c|\dd x
\end{align*}
成立,利用它即可.
\end{proof}

与 Hilbert 变换相关,有下面的结果.

\begin{theorem}\label{thm:1.13.3}
存在常数 $C$ 使
\begin{equation*}
\| Hf \|_{\mathrm{BMO}} \leq C\| f \|_{L^{\infty}} \quad (f \in L_{c}^{\infty}(\mathbb{R})).
\end{equation*}
\end{theorem}

\begin{proof}
设 $f$ 为支集是有界集合的有界函数,$Q$ 为以 $x_{0}$ 为中心,宽 $2r$ 的区间 $(x_{0} - r, x_{0} + r)$,$2Q = (x_{0} - 2r, x_{0} + 2r)$.令 $f_{1}(x) = f(x)\chi_{2Q}(x)$,$f_{2}(x) = f(x) - f_{1}(x)$.则当 $x \in Q$,$y \in (2Q)^{c}$ 时,$|x - y| \geq |x_{0} - y| - |x_{0} - x| > |x_{0} - y|/2$,故
\begin{align*}
\left|Hf(x) - \displaystyle\int_{(2Q)^{c}}\frac{f(y)}{x_{0} - y}\dd y\right| & = \left|Hf_{1}(x) + \displaystyle\int_{(2Q)^{c}}\left(\frac{1}{x - y} - \frac{1}{x_{0} - y}\right)f(y)\dd y\right| \\
& \leq |Hf_{1}(x)| + \displaystyle\int_{(2Q)^{c}}\frac{2|x - x_{0}||f(y)|}{|x_{0} - y|^{2}}\dd y \\
& \leq |Hf_{1}(x)| + 2\| f \|_{L^{\infty}}.
\end{align*}
因此,也利用 Hilbert 变换的 $L^{2}$ 有界性(定理~\ref{thm:1.3.5}),
\begin{align*}
\dfrac{1}{|Q|}\displaystyle\int_{Q}\left|Hf(x) - \int_{(2Q)^{c}}\frac{f(y)}{x_{0} - y}\dd y\right|\dd x & \leq \dfrac{1}{|Q|}\int_{Q}|Hf_{1}(x)|\dd x + 2\| f \|_{L^{\infty}} \\
& \leq \left(\dfrac{1}{|Q|}\displaystyle\int_{Q}|Hf_{1}(x)|^{2}\dd x\right)^{1/2} + 2\| f \|_{L^{\infty}}
\end{align*}
\begin{align*}
&\leq \left(\dfrac{1}{|Q|}\displaystyle\int_{\mathbb{R}}|f_{1}(x)|^{2}\dd x\right)^{1/2} + 2\| f \|_{L^{\infty}} \\
&\leq (\sqrt{2} + 2)\| f \|_{L^{\infty}}.
\end{align*}
因 $Q$ 是任意区间,由前引理~\ref{lem:1.13.2},
\begin{equation*}
\| Hf \|_{\mathrm{BMO}} \leq 2(\sqrt{2} + 2)\| f \|_{L^{\infty}}.
\end{equation*}
\end{proof}

对 Hilbert 变换稍作修正,可知它把有界函数变成 BMO 函数.

\begin{theorem}\label{thm:1.13.4}
对 $f \in L^{\infty}(\mathbb{R})$,
\begin{equation*}
\tilde{H}f(x) = \lim_{j \to \infty}\left(H\left(f\chi_{\{|x| < j\}}\right)(x) - \frac{1}{\pi}\int_{1 < |y| < j}\frac{f(y)}{-y}\dd y\right)
\end{equation*}
对几乎所有的 $x$ 存在,并在任意有界集上依 $L^{1}$ 范数收敛.进而,存在常数 $C$ 使
\begin{equation*}
\| \tilde{H}f \|_{\mathrm{BMO}} \leq C\| f \|_{L^{\infty}} \quad (f \in L^{\infty}(\mathbb{R})).
\end{equation*}
\end{theorem}

\begin{proof}
任意固定 $A > 0$.令 $j_{0}$ 为满足 $j_{0} \geq 2A$ 的最小整数.对 $j \in \mathbb{N}$,$j > j_{0}$,当 $|x| \leq A$ 时,
\begin{align*}
g_{j}(x) &:= H\left(f\chi_{\{|x| < j\}}\right)(x) - \dfrac{1}{\pi}\int_{1 < |y| < j}\frac{f(y)}{-y}\dd y \\
&= H\left(f\chi_{\{|x| < j_{0}\}}\right)(x) - \dfrac{1}{\pi}\int_{1 < |y| < j_{0}}\frac{f(y)}{-y}\dd y \\
&\quad + \dfrac{1}{\pi}\displaystyle\int_{j_{0} \leq |y| < j}\left(\frac{1}{x - y} - \frac{1}{-y}\right)f(y)\dd y,
\end{align*}
且
\begin{equation*}
\left|\int_{j \leq |y|}\left(\frac{1}{x - y} - \frac{1}{-y}\right)f(y)\dd y\right| \leq \| f \|_{L^{\infty}}\int_{j \leq |y|}\frac{2|x|}{|y|^{2}}\dd y \leq \frac{4A}{j}\| f \|_{L^{\infty}},
\end{equation*}
故当 $j \to \infty$ 时,
\begin{equation*}
g_{j}(x) = H\left(f\chi_{\{|x| < j\}}\right)(x) - \frac{1}{\pi}\int_{1 < |y| < j}\frac{f(y)}{-y}\dd y
\end{equation*}
在 $[-A, A]$ 上一致收敛到
\begin{equation*}
H\left(f\chi_{\{|x| < j_{0}\}}\right)(x) - \frac{1}{\pi}\int_{1 < |y| < j_{0}}\frac{f(y)}{-y}\dd y + \frac{1}{\pi}\int_{j_{0} \leq |y|}\left(\frac{1}{x - y} - \frac{1}{-y}\right)f(y)\dd y,
\end{equation*}
因此依 $L^{1}([-A, A])$ 范数收敛.

对区间 $Q$,$g_{j}$ 在 $Q$ 上依 $L^{1}$ 范数收敛到 $\tilde{H}f$,故也利用定理~\ref{thm:1.13.3},
\begin{align*}
\dfrac{1}{|Q|}\displaystyle\int_{Q}\left|\tilde{H}f(x) - \frac{1}{|Q|}\int_{Q}\tilde{H}f(y)\dd y\right|\dd x&= \displaystyle\lim_{j \to \infty}\frac{1}{|Q|}\int_{Q}\left|g_{j}(x) - \frac{1}{|Q|}\int_{Q}g_{j}(y)\dd y\right|\dd x  \\
&= \displaystyle\lim_{j \to \infty}\frac{1}{|Q|}\int_{Q}\left|H\left(f\chi_{\{|x| < j\}}\right)(x) - \frac{1}{|Q|}\int_{Q}H\left(f\chi_{\{|x| < j\}}\right)(y)\dd y\right|\dd x \\
&\leq \displaystyle\sup_{j}\| H\left(f\chi_{\{|x| < j\}}\right) \|_{\mathrm{BMO}} \leq C\| f \|_{L^{\infty}}.
\end{align*}
由此得
\begin{equation*}
\| \tilde{H}f \|_{\mathrm{BMO}} \leq C\| f \|_{L^{\infty}}.
\end{equation*}
\end{proof}

其次,考察 $\mathrm{BMO}(\mathbb{R})$ 与 $H^{1}(\mathbb{R})$ 的关系.

\begin{theorem}\label{thm:1.13.5}
若 $T$ 是 $H^{1}(\mathbb{R})$ 上的有界线性泛函,则存在有界函数 $g_{1}, g_{2}$ 使
\begin{equation*}
T(f) = \int_{\mathbb{R}}\left(f(x)g_{1}(x) + Hf(x)g_{2}(x)\right)\dd x \quad (f \in H^{1}(\mathbb{R})).
\end{equation*}
此时,$\| g_{1} \|_{L^{\infty}}, \| g_{2} \|_{L^{\infty}} \leq \| T \|_{\mathrm{op}}$.其中 $\| T \|_{\mathrm{op}} = \inf\{C > 0;\ |T(f)| \leq C\| f \|_{H^{1}}\}$.

反之,若 $g_{1}, g_{2} \in L^{\infty}(\mathbb{R})$,则
\begin{align*}
\left|\displaystyle\int_{\mathbb{R}}f(x)\left(g_{1}(x) + \tilde{H}g_{2}(x)\right)\dd x\right| & \leq \max\left(\| g_{1} \|_{L^{\infty}}, \| g_{2} \|_{L^{\infty}}\right)\| f \|_{H^{1}} \\
& \quad (f \in H_{00}^{1}(\mathbb{R})).
\end{align*}
\end{theorem}

\begin{proof}
在 $L^{1}(\mathbb{R}) \times L^{1}(\mathbb{R})$ 上赋予范数 $\| (f_{1}, f_{2}) \| = \| f_{1} \|_{L^{1}} + \| f_{2} \|_{L^{1}}$ 的空间记为 $L^{1}(\mathbb{R}) \oplus L^{1}(\mathbb{R})$,则此空间是 Banach 空间,且 $B := \{(f, Hf) \in L^{1}(\mathbb{R}) \times L^{1}(\mathbb{R});\ f \in H^{1}(\mathbb{R})\}$ 是 $L^{1}(\mathbb{R}) \oplus L^{1}(\mathbb{R})$ 的闭子空间.$H^{1}(\mathbb{R})$ 作为 Banach 空间可与 $B$ 视为同一.若 $|T(f)| \leq \| T \|_{\mathrm{op}}\| f \|_{H^{1}}$,则 $|T(f)| \leq \| T \|_{\mathrm{op}}\left(\| f \|_{L^{1}} + \| Hf \|_{L^{1}}\right)$,故 $T$ 是 $B$ 上的有界线性泛函.因此由 Hahn--Banach 定理可扩张为 $L^{1}(\mathbb{R}) \oplus L^{1}(\mathbb{R})$ 上的有界线性泛函.因 $L^{1}(\mathbb{R}) \oplus L^{1}(\mathbb{R})$ 的对偶空间是 $L^{\infty}(\mathbb{R}) \oplus L^{\infty}(\mathbb{R})$,故存在有界函数 $g_{1}, g_{2}$ 满足
\begin{equation*}
\begin{aligned}
T(f) &= \displaystyle\int_{\mathbb{R}}\left(f(x)g_{1}(x) + Hf(x)g_{2}(x)\right)\dd x, \\
\| g_{1} \|_{L^{\infty}}, \| g_{2} \|_{L^{\infty}} &\leq \| T \|_{\mathrm{op}}.
\end{aligned}
\end{equation*}

反之,设 $g_{1}, g_{2}$ 为有界函数,$f \in H_{00}^{1}(\mathbb{R})$.因 $f \in \mathcal{S}(\mathbb{R})$,$|f(x)| \leq C(1 + |x|)^{-2}$.另外,$\int_{\mathbb{R}}f(x)\dd x = 0$.

若能证明对任意有界函数 $g$,
\begin{equation}\tag{1.13.1}
\int_{\mathbb{R}}f(x)\tilde{H}g(x)\dd x = -\int_{\mathbb{R}}Hf(x)g(x)\dd x,
\label{eq:1.13.1}\end{equation}
则
\begin{align*}
\left|\displaystyle\int_{\mathbb{R}}f(x)\left(g_{1}(x) + \tilde{H}g_{2}(x)\right)\dd x\right| &= \left|\int_{\mathbb{R}}\left(f(x)g_{1}(x) - Hf(x)g_{2}(x)\right)\dd x\right| \\
&\leq \left(\| f \|_{L^{1}}\| g_{1} \|_{L^{\infty}} + \| Hf \|_{L^{1}}\| g_{2} \|_{L^{\infty}}\right) \leq \max\left(\| g_{1} \|_{L^{\infty}}, \| g_{2} \|_{L^{\infty}}\right)\| f \|_{H^{1}}.
\end{align*}
于是来证明 \eqref{eq:1.13.1}.令 $\chi_{a} = \chi_{\{x \in \mathbb{R};\ |x| \leq a\}}$.依次注意 $Hf \in L^{1}(\mathbb{R})$,$g \in L^{\infty}(\mathbb{R})$,$f, g\chi_{2^{k}} \in L^{2}(\mathbb{R})$,$\int_{\mathbb{R}}f(x)\dd x = 0$,
\begin{align*}
-\displaystyle\int_{\mathbb{R}}Hf(x)g(x)\dd x &= -\lim_{k \to \infty}\int_{\mathbb{R}}Hf(x)g(x)\chi_{2^{k}}(x)\dd x \\
&= \displaystyle\lim_{k \to \infty}\int_{\mathbb{R}}f(x)H(g\chi_{2^{k}})(x)\dd x \\
&= \displaystyle\lim_{k \to \infty}\int_{\mathbb{R}}f(x)\left(H(g\chi_{2^{k}})(x) + \frac{1}{\pi}\int_{1 \leq |y| \leq 2^{k}}\frac{g(y)}{y}\dd y\right)\dd x
\end{align*}
\begin{align*}
&= \displaystyle\lim_{k \to \infty}\int_{|x| \leq 2}f(x)\left(H(g\chi_{2^{k}})(x) + \frac{1}{\pi}\int_{1 \leq |y| \leq 2^{k}}\frac{g(y)}{y}\dd y\right)\dd x \\
&+ \displaystyle\lim_{k \to \infty}\sum_{j=1}^{\infty}\int_{2^{j} < |x| \leq 2^{j+1}}f(x)\left(H(g\chi_{2^{k}})(x) + \frac{1}{\pi}\int_{1 \leq |y| \leq 2^{k}}\frac{g(y)}{y}\dd y\right)\dd x.
\end{align*}
令
\begin{equation*}
a_{j}(k) = \int_{2^{j} < |x| \leq 2^{j+1}}f(x)\left(H(g\chi_{2^{k}})(x) + \frac{1}{\pi}\int_{1 \leq |y| \leq 2^{k}}\frac{g(y)}{y}\dd y\right)\dd x,
\end{equation*}
则

(i) 当 $k \leq j + 2$ 时,
\begin{align*}
a_{j}(k)&\leq \displaystyle\sup_{2^{j} < |x| \leq 2^{j+1}}|f(x)|\int_{2^{j} < |x| \leq 2^{j+1}}\left(|H(g\chi_{2^{k}})(x)| + \frac{1}{\pi}\int_{1 \leq |y| \leq 2^{k}}\frac{|g(y)|}{|y|}\dd y\right)\dd x  \\
&\leq C2^{-2j}\left\{\left(\displaystyle\int_{2^{j} < |x| \leq 2^{j+1}}\dd x\right)^{1/2}\left(\int_{2^{j} < |x| \leq 2^{j+1}}|H(g\chi_{2^{k}})(x)|^{2}\dd x\right)^{1/2}\right. \\
&\quad + \left.2^{j+2}k\dfrac{\log 2}{\pi}\| g \|_{L^{\infty}}\right\} \\
&\leq C2^{-2j}\left\{2^{(j+1)/2}\| g\chi_{2^{k}} \|_{L^{2}} + 2^{j+3}(j + 2)\dfrac{\log 2}{\pi}\| g \|_{L^{\infty}}\right\} \\
&\leq Cj2^{-j}\| g \|_{L^{\infty}}.
\end{align*}

(ii) 当 $k > j + 2$ 时,
\begin{align*}
H(g\chi_{2^{k}})(x) + \dfrac{1}{\pi}\displaystyle\int_{1 \leq |y| \leq 2^{k}}\frac{g(y)}{y}\dd y & = H(g\chi_{2^{j+2}})(x) + \dfrac{1}{\pi}\int_{1 \leq |y| \leq 2^{j+2}}\frac{g(y)}{y}\dd y \\
& \quad + \dfrac{1}{\pi}\displaystyle\int_{2^{j+2} \leq |y| \leq 2^{k}}\left(\frac{1}{x - y} + \frac{1}{y}\right)g(y)\dd y,
\end{align*}
且当 $2^{j} < |x| \leq 2^{j+1}$ 时,
\begin{align*}
\left|\displaystyle\int_{2^{j+2} \leq |y| \leq 2^{k}}\left(\frac{1}{x - y} + \frac{1}{y}\right)g(y)\dd y\right| & \leq \int_{2^{j+2} \leq |y| \leq 2^{k}}\frac{|x|}{|x - y||y|}|g(y)|\dd y \\
& \leq \displaystyle\int_{2^{j+2} \leq |y| \leq 2^{k}}\frac{2|x|}{|y|^{2}}|g(y)|\dd y \leq C\| g \|_{L^{\infty}},
\end{align*}
故把 (i) 中 $j + 2 = k$ 的情形也用于此,
\begin{equation*}
a_{j}(k) \leq Cj2^{-j}\| g \|_{L^{\infty}}.
\end{equation*}
(i),(ii) 合并,$a_{j}(k) \leq Cj2^{-j}\| g \|_{L^{\infty}}$($j = 1, 2, \ldots$),故
\begin{equation*}
\lim_{k \to \infty}\sum_{j=1}^{\infty}a_{j}(k) = \sum_{j=1}^{\infty}\lim_{k \to \infty}a_{j}(k).
\end{equation*}
因此,
\begin{align*}
-\displaystyle\int_{\mathbb{R}}Hf(x)g(x)\dd x&= \displaystyle\int_{|x| \leq 2}f(x)\left(\lim_{k \to \infty}H(g\chi_{2^{k}})(x) + \frac{1}{\pi}\int_{1 \leq |y| \leq 2^{k}}\frac{g(y)}{y}\dd y\right)\dd x  \\
&\quad + \displaystyle\sum_{j=1}^{\infty}\int_{2^{j} < |x| \leq 2^{j+1}}f(x)\lim_{k \to \infty}\left(H(g\chi_{2^{k}})(x) + \frac{1}{\pi}\int_{1 \leq |y| \leq 2^{k}}\frac{g(y)}{y}\dd y\right)\dd x \\
&= \displaystyle\int_{\mathbb{R}}f(x)\lim_{k \to \infty}\left(H(g\chi_{2^{k}})(x) + \frac{1}{\pi}\int_{1 \leq |y| \leq 2^{k}}\frac{g(y)}{y}\dd y\right)\dd x \\
&= \displaystyle\int_{\mathbb{R}}f(x)\tilde{H}g(x)\dd x.
\end{align*}
于是式~\eqref{eq:1.13.1} 得证.
\end{proof}

再进一步,$\mathrm{BMO}(\mathbb{R})$ 函数具有 $g_{1} + \tilde{H}g_{2}$($g_{1}, g_{2} \in L^{\infty}(\mathbb{R})$)的形式,但因篇幅关系,这里不加证明地列出下面的定理.

\begin{theorem}\label{thm:1.13.6}
$H^{1}(\mathbb{R})$ 的对偶空间在下面的意义下是 $\mathrm{BMO}(\mathbb{R})$.

(a) 若 $g \in \mathrm{BMO}(\mathbb{R})$,则对 $f \in H_{00}^{1}(\mathbb{R})$,
\begin{equation}\tag{1.13.2}
\ell(f) := \int_{\mathbb{R}}f(x)g(x)\dd x
\label{eq:1.13.2}\end{equation}
存在,且
\begin{equation*}
\left|\int_{\mathbb{R}}f(x)g(x)\dd x\right| \leq \| f \|_{H^{1}}\| g \|_{\mathrm{BMO}},
\end{equation*}
$\ell$ 可作为 $H^{1}(\mathbb{R})$ 上的有界线性泛函唯一扩张.

(b) 反之,若 $\ell$ 是 $H^{1}(\mathbb{R})$ 上的有界线性泛函,则存在对 $f \in H_{00}^{1}(\mathbb{R})$ 满足式~\eqref{eq:1.13.2} 的 $\mathrm{BMO}(\mathbb{R})$ 函数 $g$,且除常数函数之差外唯一确定.这个 $g$ 与定理~\ref{thm:1.13.5} 所得者除常数之差外一致.
\end{theorem}

下面不加证明地列出关于 $\mathrm{BMO}(\mathbb{R})$ 的基本事实.

\begin{lemma}\label{lem:1.13.7}
(i) 若 $f \in L^{\infty}(\mathbb{R})$,则 $\| f \|_{\mathrm{BMO}} \leq 2\| f \|_{L^{\infty}}$.

(ii) 若 $f \in \mathrm{BMO}(\mathbb{R})$,$\alpha > 1$,则
\begin{equation*}
\int_{\mathbb{R}}\frac{|f(x)|}{(1 + |x|)^{\alpha}}\dd x \leq C_{\alpha}\| f \|_{\mathrm{BMO}}.
\end{equation*}

(iii) 若 $|F(x) - F(y)| \leq C_{0}|x - y|$($x, y \in \mathbb{R}$),则对实值 BMO 函数 $f$,
\begin{equation*}
\| F(f) \|_{\mathrm{BMO}} \leq 2C_{0}\| f \|_{\mathrm{BMO}}.
\end{equation*}
由此,若 $f, g \in \mathrm{BMO}(\mathbb{R})$ 且为实值,则
\begin{align*}
\| \max(f, g) \|_{\mathrm{BMO}} &\leq 2\max\left(\| f \|_{\mathrm{BMO}}, \| g \|_{\mathrm{BMO}}\right), \\
\| \min(f, g) \|_{\mathrm{BMO}} &\leq 2\max\left(\| f \|_{\mathrm{BMO}}, \| g \|_{\mathrm{BMO}}\right).
\end{align*}
\end{lemma}

\begin{theorem}[John--Nirenberg]\label{thm:1.13.8}
若 $f \in \mathrm{BMO}(\mathbb{R})$,则对任意区间 $Q$,存在正定数 $C_{0}, C_{1}$ 使
\begin{equation*}
\frac{\left|\{x \in Q;\ |f(x) - f_{Q}| > \lambda\}\right|}{|Q|} \leq C_{0}e^{-C_{1}\lambda/\| f \|_{\mathrm{BMO}}} \quad (\lambda > 0).
\end{equation*}
它们不依赖于 $f, \lambda, Q$.
\end{theorem}

由此,

\begin{lemma}\label{lem:1.13.9}
设 $1 < p < \infty$.$f \in \mathrm{BMO}(\mathbb{R})$ 的充要条件是 $f \in L_{\mathrm{loc}}^{p}(\mathbb{R})$ 且
\begin{equation*}
\| f \|_{\mathrm{BMO}, p} := \sup_{Q:\ \text{区间}}\left(\frac{1}{|Q|}\int_{Q}|f(x) - f_{Q}|^{p}\dd x\right)^{1/p} < \infty.
\end{equation*}
此时,存在正定数 $C_{p}$ 使 $\| f \|_{\mathrm{BMO}} \leq \| f \|_{\mathrm{BMO}, p} \leq C_{p}\| f \|_{\mathrm{BMO}}$.
\end{lemma}

关于在 $\mathrm{BMO}(\mathbb{R})$ 中弱意义下稠密的好的空间,

\begin{lemma}\label{lem:1.13.10}
设 $f \in \mathrm{BMO}(\mathbb{R})$,则存在函数列 $f_{n} \in C_{c}^{\infty}(\mathbb{R})$ 满足
\begin{equation*}
\begin{aligned}
\| f_{n} \|_{\mathrm{BMO}} &\leq 2\| f \|_{\mathrm{BMO}}, \\
\lim_{n \to \infty}\displaystyle\int_{\mathbb{R}}f_{n}(x)\varphi(x)\dd x &= \int_{\mathbb{R}}f(x)\varphi(x)\dd x \quad (\varphi \in \mathcal{S}(\mathbb{R})).
\end{aligned}
\end{equation*}
\end{lemma}
